\Ch{Resource Types}{Resources}

\p Resource types are intangible types representing memory or configuration with an external interface.

\p These take the form of typed buffers, raw buffers, textures, constant buffers
and samplers.

\p Typed buffer, raw buffer, and Texture resources may be read-only or
read-write as denoted by the type of the resource, where the names of read-write
resource types are prefixed with \texttt{RW}.

\p Global resource declarations are \textit{bound} to device memory managed by
an application through a runtime implementation. A resource declaration that is
not bound is said to be \textit{unbound}. Operations on unbound resources are
evaluated as no-ops.

\p An implementation may impose restrictions on how local resource objects may
be used (\ref{AnnexB.LocalResources}).

\Sec{Typed Buffers}{Resources.TypedBuffers}

\p The typed buffer class template represents a one-dimensional resource
containing an array of a single given type. Its contents are indexed by typed
access to each element. Types may have their formats converted upon load in an
implementation-defined way.

\p The element type may be any type up to 16 bytes in size. Elements may be
padded to 16 bytes if the element type is smaller than 16 bytes.

\p All typed buffers can be read through subscript operators or Load methods.
Writable typed buffers can be written through subscript operators.

\p Where \texttt{T} can be any arithmetic type \ref{Basic.Types.Arithmetic}
or a vector with a maximum of four elements containing such an arithmetic type.
The total size of a single contained element must be less than 16 bytes.

\p Typed buffers perform format conversions on load such that the underlying
data gets converted to the destination type in an implementation-defined way.

\p The \texttt{Buffer} type's template argument \texttt{T} has a default value
of \texttt{float4}.

\begin{HLSL}
template <typename T = float4> requires hlsl::is_vector<T>::value || hlsl::is_arithmetic<T>::value
class Buffer {
  Buffer();
  Buffer(const Buffer &buf);
  Buffer &operator=(Buffer &buf);
  Buffer(const __DynamicResourceHandle &handle);

  void GetDimensions(out uint width) const;

  // Element access.
  T operator[](int index) const;
  T Load(int index) const;
  T Load(int index, out uint status) const;
};


template <typename T> requires hlsl::is_vector<T>::value || hlsl::is_arithmetic<T>::value
class RWBuffer {
  RWBuffer();
  RWBuffer(RWBuffer buf);
  RWBuffer &operator=(RWBuffer &buf);
  Buffer(const __DynamicResourceHandle &handle);

  void GetDimensions(out uint width) const;

  // Element access.
  T operator[](int index) const;
  device T &operator[](int n) const;
  T Load(int index) const;
  T Load(int index, out uint status) const;
};
\end{HLSL}

\Sub{Constructors}{Resources.TypedBuffers.Ctrs}

\p When defined at global scope, typed buffers are bound to externally-defined
backing stores using the explicit binding location if provided or the implicit
binding if not (\ref{Resources.Binding}).

\p When defined at local scope, typed buffers represent local references that
can be associated with global buffers when assigned. Use of an uninitialized
typed buffer is undefined.

\begin{HLSL}
Buffer<int> grobuf;
RWBuffer<int> grwbuf;
void main() {
  Buffer<int> lrobuf = grobuf;
  RWBuffer<int> lrwbuf = grwbuf;
}
\end{HLSL}

\p Buffers cannot be implicitly nor explicitly cast from one type to another.

\p Buffers may be created from dynamic resources. The
\texttt{\_\_DynamicResourceHandle} type in the examples above is a placeholder
for the implementation-specific type returned by subscript indexing into the
\texttt{ResourceDescriptorHeap} (\ref{Resources.Binding.Dynamic}).

\Sub{Dimensions}{Resources.TypedBuffers.Dims}

\p The size of a typed buffer can be retrieved using the \texttt{GetDimensions}
method.

\begin{HLSL}
void GetDimensions(out uint width) const;
\end{HLSL}

\p This returns the full length, in elements of the buffer through the
\texttt{width} output parameter.

\Sub{Element Access}{Resources.TypedBuffers.Access}

\p The contents of typed buffers can be retrieved using the \texttt{Load} methods
or the subscript operator.

\begin{HLSL}
T Load(uint index) const;
T operator[](uint index) const;
\end{HLSL}

\p These each return the element of the buffer's type at the given index
\texttt{index}.

\p An additional \texttt{Load} method returns the indicated element just as the
other, but also takes a second \texttt{status} parameter that returns
information about the accessed resource.

\begin{HLSL}
T Load(int index, out uint status) const;
\end{HLSL}

\p The \texttt{status} parameter returns an indication of whether all of the
retrieved value came from fully mapped parts of the resource. This parameter
must be passed to the built-in function \texttt{CheckAccessFullyMapped}
(\ref{Resources.Mapcheck}) in order to interpret it.

\p Writable buffers have an additional subscript operator that allows assignment
to an element of the buffer. It behaves as if it returns a reference to that
element as stored in device memory.

\begin{HLSL}
device T &operator[](int index);
\end{HLSL}

\begin{note}
  Accurately modeling the memory operations for writable typed buffers would
  require a proxy-object interface since the address written to isn't actually
  the device memory since it may be converted or otherwise normalized by the
  implementation.
\end{note}

\p Writes to typed buffer elements overwrite the entire element, even if only
part of the data is written. A partial write to a typed buffer is a load of the
full value, an insert of the partial elements, then a write of the full value.

\Sec{CheckAccessFullyMapped}{Resources.mapcheck}

\p The mapped status value returned by certain texture and buffer methods can be
intrepreted by a built-in function:

\begin{HLSL}
bool CheckAccessFullyMapped(uint status);
\end{HLSL}

\p This function returns true if the value was accessed from memory that was
mapped by the runtime implementation. If any part of the return value was from
unmapped memory, false is returned.
